\chapter{Aritmética: divisores e números primos}\label{ch:g9:arith}

A aritmética estuda números inteiros e como eles se dividem. É
personagens centrais são os \omterm{def:g9:arith:prime}{números primos}, os blocos de construção a partir dos quais
todo número inteiro é montado por multiplicação. O capítulo termina com o
máximo \omterm{def:g5:division:multiple}{divisor} comum, a ferramenta certa para simplificar frações uma vez
e para todos. Esta história continua, muito mais longe, no Ensino Médio
\omterm{def:g6:measure:volume}{volume} e além.

\section{Divisores e múltiplos}

\begin{definition}[Divisor, múltiplo]\label{def:g9:arith:divisor}
Deixar $a $ e $ b $ sejam inteiros positivos. Nós dizemos isso $ b $
\emph{dividir}\index{dividir} $ a $ (ou isso $ b $ é um \omterm{def:g5:division:multiple}{divisor} de $ a $, ou
isso $ a $ é um \emph{múltiplo}\index{múltiplo} de $ b $) quando $ a = b \times
k $ para algum número inteiro $ k $ --- isto é, quando a divisão de $ a $ por $ b $
folhas resto $0$.
\end{definition}

\begin{example}
O divisores de $24$ são $1, 2, 3, 4, 6, 8, 12, 24$ --- eles entram
pares cujos produto é $24$: $(1,24)$, $(2,12)$, $(3,8)$, $(4,6)$. O
múltiplos de $7$ são $7, 14, 21, 28, \dots $
\end{example}

\begin{proposition}[Regras de divisibilidade]\label{prop:g9:arith:rules}
Um número inteiro é divisível:
\begin{itemize}
  \item por $2$ quando seu último dígito é par ($0, 2, 4, 6, 8$);
  \item por $5$ quando seu último dígito é $0$ ou $5$;
  \item por $10$ quando seu último dígito é $0$;
  \item por $3$ (resp.\ $9$) quando o soma de seus dígitos é divisível por
        $3$ (resp.\ $9$);
  \item por $4$ quando seus dois últimos dígitos formam um número divisível por $4$.
\end{itemize}
\end{proposition}
\admitted

\begin{example}
$7\,215$ termina em $5$: divisível por $5$. Seu dígito soma é
$7 + 2 + 1 + 5 = 15$, divisível por $3$ mas não $9$: então $7\,215$ é
divisível por $3$, não por $9$. De fato $7\,215 = 3 \times 5 \times 481$.
\end{example}

\section{Números primos}

\begin{definition}[Número primo]\label{def:g9:arith:prime}
UM \emph{número primo}\index{número primo} é um número inteiro $\geq 2$ cujo
apenas divisores são $1$ e ele mesmo. Os primos abaixo $30$ são
\[
2,\ 3,\ 5,\ 7,\ 11,\ 13,\ 17,\ 19,\ 23,\ 29 .
\]
O número $1$ é \emph{não} primo (por convenção) e um número inteiro
$\geq 2$ que não é primo é chamado \emph{composto}.
\end{definition}

\begin{theorem}[Fatoração principal]\label{thm:g9:arith:factorization}
Cada número inteiro $\geq 2$ é um produto de números primos, e isso
a fatoração é única até a ordem dos fatores.
\end{theorem}
\admitted

\begin{method}[Fatorando um número inteiro]\label{met:g9:arith:factorization}
Divida pelo menor primo possível, repetidamente, até chegar $1$:
\begin{enumerate}
  \item tentar $2$ desde que o número seja par;
  \item então tente $3$, então $5$, então $7$, \dots\ (somente primos);
  \item pare quando o quociente é $1$; coletar os fatores com
        expoentes.
\end{enumerate}
Basta tentar números primos $ p $ com $ p^2$ não excedendo a corrente
número: se nenhum dividir isso, o próprio número é primo.
\end{method}

\begin{example}\label{ex:g9:arith:factorization}
Fator $360$, uma divisão de cada vez:
\[
360 = 2 \times 180, \quad
180 = 2 \times 90, \quad
90 = 2 \times 45, \quad
45 = 3 \times 15, \quad
15 = 3 \times 5,
\]
então
\[
360 = 2 \times 2 \times 2 \times 3 \times 3 \times 5 = 2^3 \times 3^2
\times 5 .
\]
\end{example}

\begin{omfigure}
\begin{tikzpicture}[scale=0.78]
  \node (n360) at (0,0) {$360$};
  \node (n180) at (1.2,-1) {$180$};
  \node (n90) at (2.4,-2) {$90$};
  \node (n45) at (3.6,-3) {$45$};
  \node (n15) at (4.8,-4) {$15$};
  \node (n5) at (6.0,-5) {$5$};
  \node[omThm] (p2a) at (-1.2,-1) {$2$};
  \node[omThm] (p2b) at (0,-2) {$2$};
  \node[omThm] (p2c) at (1.2,-3) {$2$};
  \node[omThm] (p3a) at (2.4,-4) {$3$};
  \node[omThm] (p3b) at (3.6,-5) {$3$};
  \draw (n360) -- (n180); \draw (n360) -- (p2a);
  \draw (n180) -- (n90); \draw (n180) -- (p2b);
  \draw (n90) -- (n45); \draw (n90) -- (p2c);
  \draw (n45) -- (n15); \draw (n45) -- (p3a);
  \draw (n15) -- (n5); \draw (n15) -- (p3b);
\end{tikzpicture}

{\small A árvore de fatores de $360$: cada etapa divide o menor primo
fator (em vermelho). Lendo as folhas vermelhas e o final $5$:
$360 = 2^3 \times 3^2 \times 5$.}
\end{omfigure}

\begin{theorem}[Euclides]\label{thm:g9:arith:euclid}
Existem infinitamente muitos números primos.
\end{theorem}

\begin{proof}
Suponha que houvesse apenas um número finito, digamos $ p_1, p_2, \dots, p_k $e
considere
\[
N = p_1 \times p_2 \times \dots \times p_k + 1 .
\]
Dividindo $ N $ por qualquer $ p_i $ folhas resto $1$, então não $ p_i $ dividir $ N $.
Mas $ N \geq 2$ tem pelo menos um primo divisor
(\cref{thm:g9:arith:factorization}) --- um primo que não está em nossa lista.
Contradição: nenhuma lista finita pode conter todos os primos.
\end{proof}

\section{O maior divisor comum}

\begin{definition}[GCD]\label{def:g9:arith:gcd}
O \emph{divisor máximo comum}\index{divisor máximo comum} de dois
inteiros positivos $ a $ e $ b $, escrito $\gcd(a, b)$, é o maior
inteiro dividindo ambos. Quando $\gcd(a, b) = 1$, os inteiros são chamados
\emph{coprime}\index{coprime}: eles não compartilham divisor exceto $1$.
\end{definition}

\begin{example}
Divisores de $18$: $1, 2, 3, 6, 9, 18$. Divisores de $24$: $1, 2, 3, 4, 6,
8, 12, 24$. Comum divisores: $1, 2, 3, 6$; então $\gcd(18, 24) = 6$. O
inteiros $15$ e $28$ são coprime.
\end{example}

\begin{proposition}[GCD de fatorações]\label{prop:g9:arith:gcdfact}
O MDC de dois inteiros é o produto dos primos que aparecem em
\emph{ambos} fatorações, cada uma tomada com o \emph{menor} dos seus dois
expoentes.
\end{proposition}
\admitted

\begin{example}\label{ex:g9:arith:gcdfact}
$360 = 2^3 \times 3^2 \times 5$ e $84 = 2^2 \times 3 \times 7$. Comum
primos: $2$ (expoentes $3$ e $2$: manter $2$) e $3$ (expoentes $2$ e
$1$: manter $1$). Então
\[
\gcd(360, 84) = 2^2 \times 3 = 12 .
\]
\end{example}

\begin{theorem}[Algoritmo euclidiano]\label{thm:g9:arith:euclidalgo}
Se $ a = bq + r $ é a divisão de $ a $ por $ b $ com resto $ r $, então
\[
\gcd(a, b) = \gcd(b, r).
\]
Repetindo as divisões até resto é $0$, o GCD de $ a $ e $ b $ é
o \emph{último diferente de zero resto}.
\end{theorem}

\begin{proof}
De $ a = bq + r $: qualquer número inteiro dividindo $ b $ e $ r $ dividir $ bq + r = a $;
e de $ r = a - bq $: qualquer número inteiro dividindo $ a $ e $ b $ dividir $ r $. Então
os pares $(a, b)$ e $(b, r)$ têm exatamente o mesmo comum divisores ---
em particular o mesmo maior. Desde restos diminuir estritamente,
o algoritmo termina, e $\gcd(x, 0) = x $ dá o último diferente de zero
resto.
\end{proof}

\begin{example}\label{ex:g9:arith:euclidalgo}
Calcular $\gcd(1071, 462)$:
\begin{align*}
1071 &= 462 \times 2 + 147, \\
462 &= 147 \times 3 + 21, \\
147 &= 21 \times 7 + 0 .
\end{align*}
O último diferente de zero resto é $21$: $\gcd(1071, 462) = 21$.
\end{example}

\begin{method}[Simplificando completamente uma fração]\label{met:g9:arith:simplify}
Para escrever $\dfrac ab $ em termos mais baixos:
\begin{enumerate}
  \item calcular $ d = \gcd(a, b)$, por exemplo\ pelo algoritmo euclidiano;
  \item dividir numerador e denominador por $ d $:
        $\dfrac ab = \dfrac{a \div d}{b \div d}$;
  \item o resultante fração é \emph{irredutível}: isso é numerador e
        denominador são coprime.
\end{enumerate}
\end{method}

\begin{example}
$\dfrac{462}{1071} = \dfrac{462 \div 21}{1071 \div 21} = \dfrac{22}{51}$,
e $\gcd(22, 51) = 1$: irredutível.
\end{example}

\section{Exercícios}

\begin{exercise}[$\star $]\label{exo:g9:arith:1}
Liste todos os divisores de $36$, de $45$, e de $17$.
\end{exercise}

\begin{exercise}[$\star $]\label{exo:g9:arith:2}
Usando o divisibilidade regras, determine se $2\,346$ é divisível por
$2$, por $3$, por $4$, por $5$, por $9$.
\end{exercise}

\begin{exercise}[$\star $]\label{exo:g9:arith:3}
Dê a fatoração primária de $72$, $150$, $210$ e $121$.
\end{exercise}

\begin{exercise}[$\star $]\label{exo:g9:arith:4}
É $101$ melhor? É $91$? É $143$? Justifique usando a regra de parada de
\cref{met:g9:arith:factorization}.
\end{exercise}

\begin{exercise}[$\star $]\label{exo:g9:arith:5}
Calcular $\gcd(48, 60)$ de duas maneiras: listando comuns divisores, e de
as fatorações primárias.
\end{exercise}

\begin{exercise}[$\star\star $]\label{exo:g9:arith:6}
Use o algoritmo euclidiano para calcular $\gcd(255, 154)$, então
$\gcd(1053, 325)$. Escreva cada linha de divisão.
\end{exercise}

\begin{exercise}[$\star\star $]\label{exo:g9:arith:7}
Faça o fração $\dfrac{588}{504}$ irredutível. (Calcule o GCD pelo
método de sua escolha e depois divida.)
\end{exercise}

\begin{exercise}[$\star\star $]\label{exo:g9:arith:8}
Uma florista tem $84$ rosas e $126$ tulipas e quer fazer idênticas
buquês, usando todas as flores, com tantos buquês quanto possível.
Quantos buquês ela pode fazer e o que cada um contém?
\end{exercise}

\begin{exercise}[$\star\star $]\label{exo:g9:arith:9}
Duas balsas saem do mesmo cais às 8h. Um parte a cada $24$ minutos,
o outro a cada $36$ minutos. A que horas eles partirão juntos?
(Procure o menor valor comum múltiplo de $24$ e $36$; fatorações
ajuda.)
\end{exercise}

\begin{exercise}[$\star\star\star $]\label{exo:g9:arith:10}
Deixar $ n $ seja um número inteiro positivo.
\begin{enumerate}
  \item Mostre isso $\gcd(n, n+1) = 1$ (números inteiros consecutivos são sempre
        coprime).
  \item Deduza que o fração $\dfrac{n}{n+1}$ é sempre irredutível.
\end{enumerate}
\end{exercise}

\section{Problema: Jarras de água, cigarras e cem armários}

\begin{problem}[{Problema de fim de semana --- o GCD decide qual
quantidades que dois jarros podem medir; os primos protegem as cigarras; e o
armários que ficam abertos são os quadrados perfeitos}]
\label{pb:g9:arith:1}
Três quebra-cabeças que parecem enigmas e são realmente aritméticos:
medindo água com jarros sem identificação (o GCD disfarçado), inseto
ciclos de vida que evoluíram para números primos, e um famoso
corredor de cem armários cujo estado final é decidido
contando divisores. Tudo funciona nas máquinas deste capítulo:
divisibilidade, fatoração primária
(\cref{thm:g9:arith:factorization}) e algoritmo de Euclides
(\cref{thm:g9:arith:euclidalgo}).

\medskip
\noindent\textbf{Parte I --- Os jarros de água.}
Você está diante de uma fonte com dois jarros sem identificação, claro. $5$ Terra
$3$ L. Movimentos permitidos: encher uma jarra até a borda, esvaziar uma jarra
completamente, despeje um jarro no outro até que a fonte esteja
vazio ou o alvo está cheio.
\begin{enumerate}
  \item Meça exatamente $1$ L. (Descreva sua sequência de movimentos
        e o conteúdo dos dois jarros depois de cada um.)
  \item Meça exatamente $4$ L --- o quebra-cabeça de um famoso
        filme de ação. (Isso pode ser feito em seis movimentos.)
  \item Quais números inteiros de litros de $1$ a $8$ você pode
        exposição (em um jarro ou dividido em ambos)? Conclua o
        list, reutilizando suas sequências.
  \item Novos jarros: $6$ Terra $4$ L. Tente medir $1$ Eu --- então
        explique por que é impossível: verifique se cada um dos três
        movimentos permitidos mantém o conteúdo de cada jarro múltiplo de
        $2$, então cada valor alcançável é par.
  \item O argumento da questão 4 funciona em geral: com jarros de $ a $
        e $ b $ litros, cada quantidade alcançável é um múltiplo de
        $\gcd(a, b)$. Calcular $\gcd(6, 4)$ e $\gcd(5, 3)$ e
        diga o que a lei prevê para cada par de jarros.
\end{enumerate}

\medskip
\noindent\textbf{Parte II --- Euclides na fonte.}
\begin{enumerate}[resume]
  \item Calcule com o algoritmo de Euclides: $\gcd(91, 65)$ e
        $\gcd(2\,026, 46)$.
  \item Explique, com suas próprias palavras, por que os valores que aparecem em
        os jarros são de Euclides restos disfarçado: com jarros
        de $13$ Terra $5$ L, encha repetidamente o jarro pequeno e
        despeje no jarro grande (esvaziando o jarro grande sempre que
        enche). Quais novos valores aparecem primeiro --- e
        compará-los com os restos no algoritmo de Euclides
        para $(13, 5)$.
  \item Deduza a resposta do campeão: com jarros de $13$ e $5$
        litros, você pode medir exatamente $1$ Eu? Justifique em um
        de acordo com a questão 5 e $\gcd(13, 5)$.
  \item Uma rápida prova de coprimalidade no estilo de
        \cref{exo:g9:arith:10}: mostre isso $\gcd(n, 2n + 1) = 1$
        para cada inteiro positivo $ n $. (O que um comum
        divisor de $ n $ e $2n + 1$ tem que dividir?)
  \item Dois ônibus saem juntos do terminal às 7h; um
        parte a cada $12$ minutos, o outro a cada $18$. Lista
        os próximos horários de partida de cada um e encontre o primeiro
        momento em que eles saem juntos novamente. Verifique neste exemplo
        a bela lei: (primeira lei comum múltiplo) $\times $
        $\gcd $ $=$ produto dos dois números --- e teste
        novamente ligado $5$ e $3$.
\end{enumerate}

\medskip
\noindent\textbf{Parte III --- Cigarras, divisores e armários.}
\begin{enumerate}[resume]
  \item Certas cigarras norte-americanas emergem apenas a cada $17$
        anos; suponha que a população de um predador atinja o pico a cada $4$
        anos. Se ambos acontecerem este ano, em quantos anos
        uma emergência em seguida coincidirá com um pico? A mesma pergunta se
        o ciclo das cigarras era $16$ anos --- com que frequência
        eles então serão massacrados? Explique em uma frase por que
        a evolução empurrou o ciclo para um \emph{melhor} comprimento.
  \item Usando a fatoração $360 = 2^3 \times 3^2 \times 5$,
        conte o divisores de $360$ sem listá-los: um
        divisor escolhe um expoente para $2$ (quatro opções:
        $0, 1, 2, 3$), um para $3$, um para $5$. Quantos
        divisores em tudo?
  \item Mostre que na fatoração de um quadrado perfeito
        $ n = m^2$, todo primo carrega um \emph{par} expoente.
        Deduza, sem calcular nenhuma raiz quadrada, que $360$ é
        não é um quadrado perfeito.
  \item Emparelhe cada um divisor $ d $ de $ n $ com seu parceiro
        $\frac{n}{d}$ (para $ n = 36$: $1 \leftrightarrow 36$,
        $2 \leftrightarrow 18$, $3 \leftrightarrow 12$,
        $4 \leftrightarrow 9$, $6 \leftrightarrow 6$). Quando é um
        divisor seu próprio parceiro? Deduza o critério: $ n $ tem um
        \emph{ímpar} número de divisores exatamente quando $ n $ é um
        quadrado perfeito. Confira $36$ e assim por diante $360$.
  \item Os cem armários. Armários $1$ para $100$ começar fechado.
        Estudante $1$ alterna todos os armários; estudante $2$ alterna
        armários $2, 4, 6, \dots $; estudante $ k $ alterna o
        múltiplos de $ k $; e assim por diante até o aluno $100$. Explique
        quais alunos tocam no armário $ n$, quantas vezes fica
        alternado e --- usando a pergunta 14 --- exatamente qual
        os armários acabam abertos. Quantos estão abertos?
\end{enumerate}
\end{problem}
